\chapter{Análisis del sistéma doméstico}
En este apartado se decidió analizar el sistema que uso laboralmente/educativamente, el cuál es un portatil MacBook air con chip M2.

% ===============================================================
% SECCIÓN 1: CONFIGURACIÓN GENERAL DEL EQUIPO
% ===============================================================
\section{Configuración General de Equipo}
\label{sec:configuracion_general}

\subsection{Especificaciones del Sistema Operativo}
\label{subsec:specs}

\textbf{Fabricante y Versión del Sistema Operativo:}
\begin{itemize}[leftmargin=*, label=\textbullet]
    \item Fabricante: Apple
    \item Versión del SO: macOS Tahoe 26.3.1
\end{itemize}

\subsection{Hardware y Componentes}
\label{subsec:hardware}

\textbf{Detalle de la Configuración del Equipo:}
\begin{itemize}[leftmargin=*, label=\textbullet]
    \item CPU: Apple M2
    \item RAM instalada: 8 GB LPDDR5
    \item Discos Duros/SSD (Capacidad y Tipo): 512 GB SSD en dos modulos de 256 GB cada uno.
    \item Otros componentes relevantes (ej. Tarjeta gráfica, etc.): GPU integrada Apple M2 con 10 núcleos,
    neural engine de 16 núcleos, cámara FaceTime HD, 2 microfonos.
\end{itemize}

% ===============================================================
% SECCIÓN 2: POLÍTICAS DE SEGURIDAD Y GESTIÓN
% ===============================================================
\section{Políticas de Seguridad y Mantenimiento}
\label{sec:politicas_seguridad}

\subsection{Gestión de Cuentas de Usuario}
\label{subsec:usuarios}

\textbf{Número total de cuentas de usuario:} 1

\textbf{Cuentas con privilegios de Administrador (o equivalente):} 1 

\begin{itemize}[leftmargin=*, label=$\diamond$]
    \item Lista de usuarios administradores: NicolasReyAlonso
\end{itemize}

\subsection{Política de Contraseñas}
\label{subsec:contraseñas}

\textbf{Descripción de la Política de Contraseñas:}
\begin{itemize}[leftmargin=*, label=$\rightarrow$]
    \item Requisitos mínimos (longitud, complejidad): Apple te permite configurar una 
    contraseña de cualquier tipo, y puntua la calidad de la contraseña, pero no impone requisitos mínimos. 
    En mi caso, tengo una contraseña de 14 caracteres mínimo con mayúsculas, minúsculas, números y símbolos.
    \item Frecuencia de cambio: No hay una frecuencia de cambio obligatoria, 
    pero apple desaconseja el cambio frecuente ya que supuestamente lleva a los usuarios a
    utilizar contraseñas simples. Este dato es interesante, ya que va en contra de la creencia común de que cambiar las contraseñas frecuentemente mejora la seguridad.
     En mi caso, no tengo una frecuencia de cambio establecida, pero suelo cambiarla cada 6 meses aproximadamente, ya que independientemente de la frecuencia, el acceso principal se realiza mediante la autenticación biométrica.
    \item Uso de autenticación multifactor (MFA): Sí, tengo configurada la autenticación multifactor utilizando la aplicación de autenticación de Apple (Apple Authenticator) y también tengo configurada la autenticación biométrica (Touch ID) para acceder al sistema.
    \item Restricciones implementadas (bloqueo de cuentas inactivas, etc.): \hrulefill 
\end{itemize}

\subsection{Política de Copias de Seguridad (Backup)}
\label{subsec:copias_seguridad}

\textbf{Estrategia de Copias de Seguridad:}
\begin{itemize}[leftmargin=*, label=$\rightarrow$]
    \item Frecuencia de las copias: Cada 3 meses, o en caso de realizar cambios significativos en el sistema, en los datos o en proyectos que no se puedan recuperar a través de github.
    \item Tipo de copias (completas, incrementales, etc.): Copias completas utilizando Time Machine de Apple.
    \item Almacenamiento (local, nube, etc.): Solo local, utilizando Time Machine de Apple con un disco duro externo.
    \item Procedimiento de recuperación en caso de fallo: En caso de fallo, se puede restaurar el sistema utilizando las copias de seguridad realizadas con Time Machine. Esto se puede hacer a través del modo de recuperación de macOS, seleccionando la opción de restaurar desde una copia de seguridad de Time Machine y siguiendo las instrucciones en pantalla para seleccionar la copia de seguridad deseada.
\end{itemize}

% ===============================================================
% SECCIÓN 3: GESTIÓN DE ACTUALIZACIONES Y EVIDENCIA
% ===============================================================
\section{Política y Gestión de Actualizaciones (Patch Management)}
\label{sec:actualizaciones}

\subsection{Requisitos y Proceso de Actualización}
\label{subsec:actualizacion_proceso}

\textbf{Fecha Límite de Mantenimiento y Seguridad:} 
\begin{itemize}[leftmargin=*, label=$\rightarrow$]
    \item Fuente de información utilizada para la fecha límite: Apple software updater.
    \item Frecuencia con la que se revisan estas fechas: Cada 3 meses, o en caso de recibir una notificación de Apple sobre una actualización crítica. 
\end{itemize}

\textbf{Descripción Detallada de la Política de Actualizaciones:}
\begin{itemize}[leftmargin=*, label=$\rightarrow$]
    \item Tipo de actualizaciones permitidas (seguridad, funcionalidad, etc.): Ninguna automática. solo permito la descarga de actualizaciones pero no la instalación automática. 
    \item Frecuencia de instalación: Tras realizar la copia de seguridad
    \item Método de instalación (automático/manual): automático a través de Apple Software Updater y de brew para las aplicaciones instaladas a través de este gestor de paquetes.
    \item Procedimiento de prueba antes de la implementación: \hrulefill 
\end{itemize}

\subsection{Evidencia Gráfica}
\label{subsec:evidencia_pantalla}

\textbf{Captura de Pantalla del Nivel de Actualización del Sistema Operativo:}
\begin{figure}[h]
    \centering
    % Reemplaza 'nombre_captura.png' con el nombre real de tu archivo
    % \includegraphics[width=0.8\textwidth]{nombre_captura_actualizacion.png} 
    \caption{Captura de pantalla del nivel de actualización y estado de seguridad del Sistema Operativo.}
    \label{fig:actualizacion}
\end{figure}
